\documentclass[11pt]{article}
% Source-PDF: 图论 GRAPH THEORY/2-SAT信息学奥林匹克论文 - 赵爽.pdf
\usepackage{oipl-vn}

\title{Phân tích sơ lược lời giải 2-SAT}
\author{Zhao Shuang}
\date{}

\begin{document}
\maketitle

\origtitle{2-SAT problem analysis}
\sourcepath{graph-theory/2-SAT-informatics-olympiad-paper-zhao-shuang.pdf}

\section{Cơ sở lý thuyết SAT}

Đặt $B=\{b_1,b_2,\ldots,b_n\}$ là một tập hữu hạn các biến Boolean, và
\[
  \hat B=\{b_1,b_2,\ldots,b_n,\neg b_1,\neg b_2,\ldots,\neg b_n\}.
\]
Nếu $B'$ là một tập con không rỗng của $\hat B$, định nghĩa
\[
  \bigvee B'=\bigvee_{b\in B'} b.
\]
Với các tập đã cho $B'_1,B'_2,\ldots,B'_m\in \hat B$,\footnote{Trong phần sau,
$X_1,\ldots,X_n$ được viết gọn thành $\{X_n\}$.} bài toán tìm một phép gán cho
$B$ sao cho
\[
  \left(\bigvee B'_1\right)\land
  \left(\bigvee B'_2\right)\land \cdots \land
  \left(\bigvee B'_m\right)=1
\]
được gọi là bài toán \term{thỏa mãn được} (Satisfiability), viết tắt là SAT.

Đặc biệt, với một họ các tập $\{B'_m\}$ đã cho, nếu
\[
  \max_{i=1,2,\ldots,m} |B'_i|=k,
\]
ta gọi bài toán này là $k$-SAT. Có thể chứng minh rằng khi $k>2$, $k$-SAT là
NP-đầy đủ. Phần dưới đây xét trường hợp $k=2$.

\section{2-SAT}

Trong 2-SAT, mỗi $B'_i$ chỉ có hai dạng: một biến Boolean đơn lẻ
$x\in \hat B$, hoặc phép hoặc của hai biến Boolean:
\[
  x\lor y,\qquad x,y\in \hat B.
\]
Để tiện trình bày, trước hết ta phân tích trường hợp chỉ tồn tại dạng thứ hai.

Ta có thể dùng một đồ thị có hướng $G$. Đồ thị $G$ có $2n$ đỉnh, tương ứng với
$2n$ phần tử của $\hat B$. Bài toán được chuyển thành việc chọn $n$ đỉnh trong
$G$ sao cho thỏa mãn điều kiện của 2-SAT; tất nhiên hai đỉnh đại diện cho
$b_i$ và $\neg b_i$ không được đồng thời được chọn. Sau đây xét mệnh đề
$B'_i=x\lor y$ tương ứng với gì trên đồ thị.

Rõ ràng
\[
  x\lor y=\neg(\neg x\land \neg y).
\]
Điều này có nghĩa là nếu ta chọn $\neg x$ thì bắt buộc phải chọn $y$; tương tự,
nếu chọn $\neg y$ thì bắt buộc phải chọn $x$. Vì vậy, với $B'_i=x\lor y$, ta
thêm vào $G$ hai cung $(\neg x,y)$ và $(\neg y,x)$.\footnote{Ở đây
$x,y,\neg x,\neg y$ đều chỉ các đỉnh trong $G$ đại diện cho chúng. Các phần
sau cũng dùng quy ước này.}

Tiếp theo, ta tìm tất cả các \term{thành phần liên thông mạnh} (strongly
connected components) của $G$. Nếu chọn bất kỳ một đỉnh trong một thành phần
liên thông mạnh, thì mọi đỉnh khác trong thành phần đó cũng bắt buộc phải
được chọn. Nếu $b_j$ và $\neg b_j$ nằm trong cùng một thành phần liên thông
mạnh, ta gặp mâu thuẫn và công thức 2-SAT vô nghiệm.

Nếu không có mâu thuẫn, ta gộp các đỉnh và cung thuộc cùng một thành phần liên
thông mạnh thành một đỉnh, thu được đồ thị có hướng mới $G'$. Sau đó đảo
chiều tất cả các cung của $G'$ để thu được $G''$.

Bây giờ xét $G''$. Vì đã gộp các thành phần liên thông mạnh thành đỉnh, $G''$
không có chu trình, nghĩa là $G''$ là một DAG và có thứ tự topo. Ta đặt tất cả
các đỉnh của $G''$ ở trạng thái "chưa tô màu", rồi lặp lại các thao tác sau
theo thứ tự topo:

\begin{enumerate}
  \item Chọn đỉnh chưa tô màu đầu tiên $x$, và tô $x$ màu đỏ.
  \item Tô màu xanh cho mọi đỉnh $y$ mâu thuẫn với $x$, cùng toàn bộ hậu duệ
  của $y$. Hai đỉnh $x$ và $y$ được gọi là mâu thuẫn nếu tồn tại
  $b_j,\neg b_j\in \hat B$ sao cho $b_j$ thuộc thành phần liên thông mạnh mà
  $x$ đại diện, còn $\neg b_j$ thuộc thành phần liên thông mạnh mà $y$ đại
  diện.
  \item Lặp lại hai bước trên cho đến khi không còn đỉnh chưa tô màu. Khi đó,
  tập các đỉnh được tô màu đỏ trong $G''$, ánh xạ ngược về tập đỉnh tương ứng
  trong $G$, cho ta một nghiệm của 2-SAT.
\end{enumerate}

Các thao tác trên có thể sinh mâu thuẫn hay không? Câu trả lời là không.
Trước hết, giả sử ta chọn một đỉnh chưa tô màu $p$ trong $G''$. Mọi đỉnh mâu
thuẫn với $p$, cùng mọi hậu duệ của chúng, đều sẽ được tô màu xanh. Vì thế khi
tô $p$ màu đỏ, không có đỉnh nào mâu thuẫn với $p$ cũng đã có màu đỏ.

Đồng thời, do ta duyệt theo thứ tự topo và mỗi khi tô xanh một đỉnh, ta lập
tức tô xanh tất cả hậu duệ của nó. Nói cách khác, nếu một đỉnh $p$ không thể
được chọn, thì mọi đỉnh trực tiếp hoặc gián tiếp thỏa điều kiện "nếu chọn $q$
thì bắt buộc phải chọn $p$" cũng sẽ được tô màu xanh. Như vậy trong $G''$
không tồn tại cung $(x,y)$ với $x$ màu đỏ và $y$ màu xanh. Kết luận thu được
không thể trái với điều kiện tương ứng với bất kỳ $B'_i$ nào.

Từ hai nhận xét trên, ta có thể chứng minh rằng các thao tác tô màu không sinh
ra mâu thuẫn.

Tiếp theo chứng minh nghiệm tương ứng với $G''$ thật sự là nghiệm của công
thức 2-SAT.

\subsection*{1. Nghiệm không đồng thời chọn $b_j$ và $\neg b_j$}

Với $b_j$ và $\neg b_j$, chúng chắc chắn nằm trong hai thành phần liên thông
mạnh khác nhau của $G$; nếu không, bài toán đã bị kết luận là vô nghiệm từ
trước. Trong $G''$, hai thành phần này được đại diện bởi hai đỉnh khác nhau,
gọi là $p$ và $q$.

Rõ ràng $p$ và $q$ là hai đỉnh mâu thuẫn nhau. Vì quá trình tô màu ở trên
không sinh mâu thuẫn, $p$ và $q$ không thể đồng thời được tô màu đỏ. Do đó
nghiệm thu được không thể đồng thời chọn $b_j$ và $\neg b_j$.

\subsection*{2. Với mọi cặp $b_j,\neg b_j\in\hat B$, nghiệm chọn ít nhất một trong hai}

Lưu ý rằng với mỗi cung $(x,y)$ trong $G$, luôn tồn tại một cung "đối ngẫu"
$(\neg y,\neg x)$. Điều này xuất phát từ dạng của $B'_i$. Hệ quả là nếu $x$ và
$y$ cùng nằm trong một thành phần liên thông mạnh của $G$, thì $\neg x$ và
$\neg y$ cũng nằm trong một thành phần liên thông mạnh; nếu có đường đi từ $x$
đến $y$, thì cũng có đường đi từ $\neg y$ đến $\neg x$.

Giả sử ngược lại rằng cả $b_j$ lẫn $\neg b_j$ đều không được chọn, tức hai
thành phần liên thông mạnh $p$ và $q$ chứa chúng đều bị tô màu xanh. Ta xét hai
trường hợp.

\paragraph{Trường hợp 1: $p$ và $q$ bị tô xanh trong cùng một lần tô màu.}
Không mất tính tổng quát, giả sử sau khi tô đỏ đỉnh $r$, ta đồng thời tô xanh
$p$ và $q$. Trong $G''$, giữa $p$ và $q$ không có đường đi. Nếu có đường đi
giữa $p$ và $q$, thì do cấu trúc của $G$, sẽ cũng có đường đi theo chiều ngược
giữa $q$ và $p$, mâu thuẫn với tính chất DAG của $G''$.

Nếu giữa $p$ và $q$ không có đường đi, thì tồn tại $p',q'\in G''$ sao cho
$p'$ và $q'$ mâu thuẫn trực tiếp với $r$, $p$ là hậu duệ của $p'$, và $q$ là
hậu duệ của $q'$. Khi đó tồn tại
\[
  G(x)\in G''(r),\quad G(\neg x)\in G''(p'),
\]
và
\[
  G(y)\in G''(r),\quad G(\neg y)\in G''(q').
\]
Nhưng theo cấu trúc của $G$, $\neg x$ và $\neg y$ phải nằm trong cùng một
thành phần liên thông mạnh. Điều này mâu thuẫn với
$G(\neg x)\in G''(p')$ và $G(\neg y)\in G''(q')$.\footnote{Ở đây
$G(x)\in G''(r)$ có nghĩa là đỉnh $x$ của $G$ nằm trong thành phần liên thông
mạnh được đại diện bởi đỉnh $r$ của $G''$. Các ký hiệu sau cũng được hiểu như
vậy.}

Vì thế $p$ và $q$ không thể bị tô màu trong cùng một lần tô màu.

\paragraph{Trường hợp 2: $p$ và $q$ bị tô xanh ở hai thời điểm khác nhau.}
Không mất tính tổng quát, giả sử $q$ bị tô màu sau; cụ thể, khi tô đỏ $r$, do
$q$ mâu thuẫn với $r$ nên $q$ bị tô xanh.

Nếu $q$ và $r$ mâu thuẫn trực tiếp, tức tồn tại
\[
  G(x)\in G''(r),\quad G(\neg x)\in G''(q),
\]
lại có $G(\neg b_j)\in G''(q)$. Theo cấu trúc của $G$, suy ra
$G(b_j)\in G''(r)$, nên $r=p$. Điều này mâu thuẫn với việc $p$ đã bị tô xanh.

Nếu $q$ và $r$ mâu thuẫn gián tiếp, tức tồn tại
\[
  G(x)\in G''(r),\quad G(\neg x)\in G''(r'),
\]
và $q$ là tổ tiên của $r'$. Theo cấu trúc của $G$, $p$ tất yếu là tổ tiên của
$r$. Nhưng $p$ đã bị tô màu xanh, nên $r$ không thể được tô màu đỏ, mâu thuẫn.

Kết hợp hai trường hợp, mệnh đề được chứng minh.

\section{Độ phức tạp}

Qua chứng minh trên, thuật toán là đúng. Việc tìm các thành phần liên thông
mạnh của đồ thị có hướng và sắp xếp topo đều có độ phức tạp $O(e)$; thao tác
tô màu cũng có độ phức tạp $O(e)$. Vì vậy, toàn bộ thuật toán có độ phức tạp
$O(e)$.

\end{document}
